\section{Free monoids and free groups as test diagrams}
\label{sec:D}

All monoid colimits in this section are taken in the category of all
monoids, with unit-preserving homomorphisms. We write
$T_i^A=L_i\colim_A:\Ab^A\to\Ab$.

\begin{claim}{D.1}{Proved}{The homology and fundamental group in Lead 3}
\label{D:homology}
Let $A$ be a nonempty connected total small category, and let
$M:A\to\mathrm{FreeMon}$ be any functor, allowing arbitrary word
substitutions. Put $G(a)=M(a)^{\mathrm{gp}}$ and
$\Gamma=\colim_A^{\mathrm{Grp}}G$. For every integer $k\geq0$
there are natural isomorphisms
\[
 H_{k+1}\left(B\int_A M;\Z\right)\cong T_k^A(M_{\mathrm{ab}}),
 \qquad \pi_1\left(B\int_A M\right)\cong\Gamma.
 \tag{D.1.1}
\]
The fundamental-group isomorphism is the one induced by the
canonical functor $\int_A M\to\mathbf B\Gamma$, after choosing
its sole target object as basepoint and any source object. The
same assertions hold for a diagram of free groups.
\end{claim}
\begin{proof}
For a free monoid $F(S)$ the canonical map to its free group
$F^{\mathrm{gp}}(S)$ induces a weak equivalence of classifying
spaces. The universal covering category has objects the reduced
free-group words and arrows right multiplication by positive words.
Its generator graph is the free-group Cayley tree. Composition
triangles collapse each positive word to its generator path, so
its realization has the type of that tree and is contractible.
This verifies PRETOT F.9's free-monoid input for the specified
map; naturality for arbitrary homomorphisms follows from
naturality of group completion, without choosing compatible graph
bases. The same result is stated and proved in Lenz,
Proposition~4.4.1 and Corollary~4.4.6 \cite{CDLenz2011}.
Thus the natural homology functors are $H_0(BM)=\Z$,
$H_1(BM)=M_{\mathrm{ab}}$, and $H_q(BM)=0$ for $q\geq2$.

The bisimplicial bar model of the Grothendieck construction gives
\[
 E^2_{p,q}=T_p^A\bigl(H_q(BM;\Z)\bigr)
       \Longrightarrow H_{p+q}\left(B\int_A M;\Z\right).
\]
Since $A$ is total, its constant singleton action makes $BA$
weakly contractible, so $T_p^A(\Z)=0$ for $p>0$.
Only the row $q=1$ and the term $(0,0)$ remain. Every total degree
has a finite filtration and one possible nonzero graded term,
proving the homology formula even when $A$ and all free bases are
infinite.

For the fundamental group, the groupoid completion of $\int_A M$
has generators the groups $G(a)$ and one transport arrow for
each arrow of $A$. Its relations are the composition relations
of $A$ and conjugation by transport according to $G$.
Choose a maximal tree in the connected underlying graph of $A$.
After contracting that tree, all remaining transport loops vanish
because $BA$ is simply connected. The remaining presentation is
exactly the colimit presentation of the groups $G(a)$. The
canonical functor in the statement sends every transport arrow
to the identity and every $M(a)$-element to its colimit image,
so this calculation identifies its induced map, not merely an
abstractly isomorphic group.
\end{proof}

\begin{claim}{D.2}{Proved}{Exact indices and free-group diagrams}
\label{D:exact}
If $A$ is nonempty connected and exact, every diagram
$M:A\to\mathrm{FreeMon}$ satisfies
\[
 B\int_A M\simeq K(\Gamma,1).
\]
Consequently, for $C=\colim_A^{\mathrm{Mon}}M$,
\[
 q_M\in W_\infty
 \quad\Longleftrightarrow\quad
 BC\longrightarrow B(C^{\mathrm{gp}})\text{ is a weak equivalence}.
 \tag{D.2.1}
\]
Every diagram of free groups over such an $A$, with colimit in
groups, passes the comparison test. In particular free-group
diagrams do not characterize filteredness: the exact nonfiltered
category $\Delta_{\mathrm{inj}}$ passes all of them.
\end{claim}
\begin{proof}
We check the exactness-to-totality input in R6 at the step needed.
For an $A$-set $S$, direct sum over its fibers is an exact functor
$\Sigma_S:\Ab^{\El_A S}\to\Ab^A$, and
$\colim_{\El_A S}=\colim_A\Sigma_S$ by the cocone universal
property. Thus every connected action component has exact abelian
colimits and is acyclic. Applying this again to covering actions
associated with cyclic subgroups shows that its fundamental group
is trivial; a nontrivial cyclic subgroup would have nonzero
abelianization. Hurewicz then makes every component weakly
contractible. This is the action criterion for totality, so D.1
applies.

Let $X=B\int_A M$ and use D.1's canonical $X\to B\Gamma$.
Pull the universal covering $E\Gamma\to B\Gamma$ back to $X$.
The resulting cover $\widetilde X$ is connected and simply
connected because the map induces an isomorphism on $\pi_1$.
At the bar-model level it is the homotopy colimit over $A$ of
the corresponding $\Gamma$-coverings $X_a$ of $BM(a)$:
in each simplex the initial lift is a $\Gamma$-label, and all
faces and transition maps carry this label by the stipulated
monoid homomorphism. Each $X_a$ is a possibly disconnected
cover of a graph up to weak equivalence, hence
$H_q(X_a;\Z)=0$ for $q\geq2$.
The homology spectral sequence for these coverings has
$E^2_{p,q}=T_p^A H_q(X_a;\Z)$; exactness kills every $p>0$.
It follows that $H_j(\widetilde X;\Z)=0$ for $j\geq2$.
Its $H_1$ also vanishes by simple connectivity, so
$\widetilde X$ is weakly contractible. This proves the first
assertion without any injectivity assumption on the transition
maps or on $G(a)\to\Gamma$.

Group completion is left adjoint to inclusion of groups into
monoids, and therefore $C^{\mathrm{gp}}=\Gamma$.
The composite $X\xrightarrow{Bq_M}BC\to B\Gamma$ is the
canonical weak equivalence just proved. Two out of three gives
(D.2.1). For free-group diagrams the target itself is $B\Gamma$.
The known exactness of $\Delta_{\mathrm{inj}}$ is
Isbell--Mitchell, Theorem~3 \cite{IsbellMitchell1973}, as verified
in the supplied TOTEX E.2; its two point inclusions cannot be
post-equalized by an injection, so it is not filtered.
\end{proof}

\begin{claim}{D.3}{Proved}{The interval example in Lead 3}
\label{D:interval}
For $[n]=\{0,\ldots,n\}$ let $M([n])$ be the free monoid
on the consecutive edges of
$\{-\infty\}< [n] <\{+\infty\}$. An increasing injection,
extended by the two endpoints, sends each generator to the
ordered product of the intervening edges. This is a functor
$\Delta_{\mathrm{inj}}\to\mathrm{FreeMon}$, and its colimit is
\[
 C=\langle x,y,z\mid xy=x,\ y^2=y,\ yz=z\rangle.
 \tag{D.3.1}
\]
Its group completion is the free group $F(x,z)$ and
$H_2(BC;\Z)=0$. Thus this example does not separate the comparison
in degree two. The stronger theorem D.6 below gives
$BC\simeq K(F(x,z),1)$ and proves that its entire comparison passes.
\end{claim}
\begin{proof}
Subdividing an interval and then subdividing its pieces gives
exactly the same ordered word as the composite injection; this
proves functoriality. Every first edge is the image of the first
edge at $[0]$, using the injection selecting the first vertex.
Every last edge is similarly represented at $[0]$. Every interior
edge is the image of the sole interior edge at $[1]$, using its
two endpoints. Thus $x,y,z$ generate the colimit. The two maps
$[0]\to[1]$ give $x=xy$ and $z=yz$; the injection
$[1]\to[2]$ skipping the middle point gives $y=y^2$.
Conversely the word assigned to any first, interior, or last
edge is respectively $xy^a$, $y^{a+1}$, or $y^az$ for some
$a\geq0$, and reduces to $x,y,z$ using these three relations.
They therefore define a cocone to (D.3.1). The two resulting
maps are inverse on every original generator, proving the
presentation, with no unverified absence-of-further-relations step.

In the group completion the first relation gives $y=1$; the
other two then impose no relation on $x,z$.
For monoid homology orient the three rules as
$xy\to x$, $yy\to y$, $yz\to z$. Each decreases length.
Their only overlapping critical words are
$xyy,yyy,yyz,xyz$; both reductions have the same normal form
in each case. Disjoint reductions commute. Induction on word
length therefore gives convergence. The coherent completion
by these four critical confluences supplies a free resolution
through degree three, by the proof of Squier's theorem in
\cite[Theorem~7.3.5 and Proposition~9.3.2]{CDPolygraphs2025}.
For completeness the low-degree exactness used here follows
by choosing the leftmost normalizing path: induction compares
any competing first reduction with it by one critical or
commuting square, then treats the shorter resulting words.
This contracts the associated relation module and its module
of identities among relations, precisely the first three stages
of that resolution.

After tensoring with the trivial module $\Z$, let
$r_1,r_2,r_3$ denote the three rules. Each has boundary $y$.
The four confluence boundaries are
\[
 r_1-r_2,\quad 0,\quad r_2-r_3,\quad r_1-r_3.
\]
They generate the entire kernel of
$\Z^3\to\Z\{x,y,z\}$, $(a,b,c)\mapsto(a+b+c)y$.
Thus $H_2(BC)=0$ and $H_1(BC)=\Z^2$.
The program \texttt{CD\_finite\_checks.py} checks all critical
pairs, the cocone on every injection through $[5]$, and the
integral matrices; a unit minor verifies that the image lattice
is primitive. Convergence and the universal-property proof above
supply the statements for all lengths and all simplex dimensions.
\end{proof}

\begin{claim}{D.4}{Proved}{The free-group obstruction over the total category $\mathcal C_\delta$}
\label{D:predimension}
Let $A=\mathcal C_\delta$ be the category in R17. Let $V$ be
the one-point object and $W$ the ordered three-point one-edge
object, with point maps $v_i:V\to W$. The presentations
\[
 G(X)=\left\langle u_x,v_x\ (x:V\to X)\ \middle|\
 u_{ev_2}u_{ev_1}^{-1}v_{ev_3}v_{ev_1}^{-1}=1
 \ (e:W\to X)\right\rangle
 \tag{D.4.1}
\]
define a functor $G:A\to\mathrm{Grp}$ by postcomposition of
the indices. Every value is a finitely generated free group;
every transition is a free-factor inclusion. Its abelianization
is $Q=\operatorname{coker}(P_W\to P_V^2)$ with
\[
 D[e]=([ev_2]-[ev_1],\ [ev_3]-[ev_1]),
 \qquad T_1^A(Q)=\Z.
\]
Furthermore
\[
 \colim_A^{\mathrm{Grp}}G=F(u,v),\qquad
 H_2\left(B\int_A G;\Z\right)=\Z.
\]
The specified comparison to the group colimit is therefore not
a weak equivalence.
\end{claim}
\begin{proof}
All arrows of $A$ are increasing induced embeddings with strictly
closed image. An arrow sends every actual point and actual edge
embedding to another, and sends every displayed relator to the
corresponding relator. This proves functoriality of the actual
presentations.

Consider any nonempty family of $r$ remaining relators at $X$.
Its $W$-arrows have distinct hyperedge images, since an increasing
embedding of $W$ is determined by its image. If their union has
$v$ vertices, positivity gives $2v>3r$. The total number of
incidences is $3r$, so some vertex has degree exactly one in
this family. If it occupies position $1$, $2$, or $3$ in its
unique edge, select respectively $u_{ev_1}$, $u_{ev_2}$, or
$v_{ev_3}$. That generator occurs exactly once, with exponent
$\pm1$, in the unique relator, and in no other remaining relator.
Tietze elimination removes it and that relator while leaving
all other relators unchanged. Induction removes every relator,
leaving a finite free presentation.

For an arrow $Y\to X$, identify $Y$ with its closed induced
image. An actual point or edge embedding of $X$ supported in
$Y$ is an actual embedding in $Y$, by restriction of closedness;
the converse is transitivity. For any nonempty family of $r$
new edges let $U$ be $Y$ together with their vertices.
Strict closedness gives
\[
 2|U\setminus Y|>
 3|E_X(U)\setminus E_X(Y)|\geq3r.
\]
Some new vertex therefore occurs in just one of these edges.
The preceding elimination can always use a generator outside
$Y$. It never removes or changes an old relator. After all new
relations have been eliminated, the presentation is the old
presentation at $Y$ with freely adjoined generators. This proves
that the specified map, not just an abstract map between free
groups of the right ranks, is a free-factor inclusion.

Abelianizing (D.4.1) gives exactly $Q$. The same private-vertex
argument makes $D:P_W\to P_V^2$ injective: a nonzero finite
linear combination of columns has a coordinate supported in
one edge only and hence cannot vanish. Thus
$0\to P_W\to P_V^2\to Q\to0$ is a projective resolution.
The colimits of both representables are $\Z$ and the colimit
map $\Z\to\Z^2$ is zero. The derived-colimit exact sequence
gives $T_1Q=\Z$, with generator the class corresponding to
$1\in\colim P_W=\Z$.

The two sets of generators in (D.4.1) are representable actions.
Each has singleton colimit. Since the free-group functor preserves
colimits, the group colimit has two generators $u,v$; every
relator becomes $uu^{-1}vv^{-1}=1$. Hence the canonical colimit
is $F(u,v)$. Applying D.1 to the known total category $A$ gives
$H_2(B\int_A G)=T_1Q=\Z$, while $BF(u,v)$ is a graph and
has $H_2=0$.

We verified the totality input of TOTAL\_EXACT\_2 A.30--A.34
in the attached stabilizer and action-reduction proofs. The
alternating amalgam identifies only neighboring $B$- and $C$-copies;
nonfixed vertices have visibility at most three. Cutting a
seam loses at most $2\max(|B|,|C|)$ in predimension, whereas
covering all seams contributes at least the number of seams
divided by three. Choosing sufficiently many seams proves
closedness of all copies. The displayed interpolation preserves
the order and agrees on $B\cap C$, yielding the algebraic identity
$\langle G_B,G_C\rangle=G_{B\cap C}$. Minimal finite supports
and rigidity of finite linear orders then identify each open
subgroup with $G_F$.
For arbitrary actions, quotienting pairs eventually made equal
is functorial by condition I and gives injective transitions.
Both comparison commas in the proof are nonempty cosieves of
slices, contracted by functorial ordered free amalgamation with
a fixed slice object. Saturation by the directed union over
finite closed substructures then identifies each orbit action
with a representable $A(F,-)$, or the constant singleton if
$F$ is empty. These steps apply to all actions; no finite-action
or injective-action restriction remains. This confirms exactly
the totality hypothesis needed for the final use of D.1.
The independent program \texttt{CD\_freefactor\_checks.py}
checks all 172 positive ordered hypergraphs through five vertices,
all 775 closed-subobject embeddings, and every actual-edge family
in each; its 496 nonempty relative families all have a new private
vertex. Every relation matrix also has full row rank. These bounded
checks supplement the general counting and elimination proof.
\end{proof}

\begin{claim}{D.5}{Proved}{What the earlier free-factor examples do and do not change}
\label{D:scope}
The free-group construction of TOTAL\_EXACT\_2 C.8 over $K^*$,
and its finite free-factor example C.9, do not refute D.2.
The former refutes an extension of D.2 from exact to arbitrary
quasi-filtered indices. D.4 refutes that extension even for
total indices. The latter refutes deducing asphericity of a
simplicial nerve colimit solely from free values, free-factor
maps, and a free group colimit.
\end{claim}
\begin{proof}
We read the proof of C.8 in \texttt{akr\_links.tex}. Its absolute
and relative private-vertex eliminations are exactly the typed
group-presentation argument above; its derived-colimit calculation
is $T_1Q=\Z$. In particular that index cannot be exact, so it
does not satisfy D.2's hypothesis.

For C.9 the index is the parallel-pair category
$0\rightrightarrows1$. Its maps are
\[
 F(a,c)\rightrightarrows F(a,c,b),\qquad
 f(a)=a,\ f(c)=c,\quad g(a)=ab,\ g(c)=cb^2.
\]
The second is the restriction of the automorphism fixing $b$
and sending $a\mapsto ab$, $c\mapsto cb^2$; its inverse uses
$b^{-1},b^{-2}$. The two image subgroups intersect trivially:
in a nonempty reduced word, a negative first letter leaves a
nonzero initial $b$-block and a positive last letter leaves a
nonzero final one. In the remaining case a positive-to-negative
transition involves distinct generators, leaving exponent
$1-2$ or $2-1$. Intervening $a,c$-blocks cannot cancel because
the original word is reduced. Thus the nerve images intersect
only in their vertex, and the coequalizer is a homotopy pushout
of their wedge inclusion. Mayer--Vietoris gives
\[
 H_2\bigl(\operatorname{coeq}(Nf,Ng);\Z\bigr)
 =\ker\bigl(\Z^2\xrightarrow{(-1,-2)}\Z\bigr)
 =\Z(2,-1).
\]
The group-colimit relation $a=ab$ kills $b$ and makes the
second relation redundant, so that colimit is $F(a,c)$.
The index itself has $H_1=\Z$ and is not total or exact.
This checks the precise scope of both earlier results.
The same independent program verifies the two Nielsen inverses on
each signed basis element, the primitive integral kernel just
displayed, and the group-colimit elimination. It also checks the
surviving-$b$ assertion on all 39,364 nonempty reduced words through
length nine. The reduced-word argument above proves that assertion
without a length bound.
\end{proof}

\begin{claim}{D.6}{Proved}{Every exact index passes every free-monoid diagram}
\label{D:exact-monoids}
Let $A$ be a nonempty connected small category for which
$\colim_A:\Ab^A\to\Ab$ is exact. For every diagram
$M:A\to\mathrm{FreeMon}$, allowing all monoid homomorphisms,
the canonical comparison
\[
 q_M:\int_A M\longrightarrow
 \mathbf B\!\left(\colim_A^{\mathrm{Mon}}M\right)
\]
is homotopy final: its unique target comma category is weakly
contractible. In particular $q_M\in W_\infty$, and every such
monoid colimit has aspherical classifying space.
\end{claim}
\begin{proof}
Write $C=\colim_A^{\mathrm{Mon}}M$ and
$j_a:M(a)\to C$ for its cocone. Define a diagram $D:A\to\Cat$
whose object set is the same set $C$ at every stage. Its arrows
in stage $a$ are
\[
 \Hom_{D(a)}(c,c')=
 \{m\in M(a):j_a(m)c=c'\}.
 \tag{D.6.1}
\]
Composition is monoid multiplication. A transition $f:a\to b$
is the identity on objects and sends the arrow $m$ to $M(f)m$.
These formulas identify the comma category
$*\backslash q_M$ with $\int_A D$ exactly, not just up to
homotopy.

There are three points to check.

\emph{First, each $D(a)$ is a free category on a directed graph.}
Choose a free basis $S_a$ of $M(a)$. The graph has vertices $C$
and, for each $c\in C$ and $s\in S_a$, one edge
$c\to j_a(s)c$. A word in $S_a$ specifies its unique directed
path starting at $c$, and equality of words is precisely equality
of the resulting arrows in (D.6.1). Thus this graph freely
presents $D(a)$. No compatible choice of bases is needed:
transition functors may send an edge to an arbitrary finite
path, including a constant path. For completeness the required homology assertion has a short
resolution proof. If $F$ is the free category on a graph with
vertices $V$ and edges $e:v\to w$, put
$P_v=\Z[F(v,-)]$. Then
\[
 0\longrightarrow\bigoplus_{e:v\to w}P_w
 \xrightarrow{\partial}\bigoplus_{v\in V}P_v
 \longrightarrow\underline\Z\longrightarrow0,
 \quad \partial(p,e)=[pe]-[p],
\]
is exact. At an object $x$, its augmented basis graph is the
rooted tree of all finite paths ending at $x$: removing the first
edge gives the unique predecessor of each nonempty path. Thus
it is a projective resolution of length one, for arbitrary
infinite graphs as well. Consequently
\[
 H_q(BD(a);\Z)=0\qquad(q\geq2).
 \tag{D.6.2}
\]

\emph{Second, the strict categorical colimit is the regular
translation category of $C$.} More precisely,
\[
 \colim_A^{\Cat}D\cong D_C,
 \qquad
 \Hom_{D_C}(c,c')=\{u\in C:uc=c'\}.
 \tag{D.6.3}
\]
Here is a presentation proof, which checks the possible
nonfiltered-colimit issue. Present $C$ by the disjoint union of
the bases $S_a$, with the substitution relations
$s=M(f)(s)$ for all $f:a\to b$ and $s\in S_a$.
Present the categorical colimit by the graph edges $(c,s)$
above, imposing the same substitutions as path relations
\emph{at every initial vertex $c\in C$}. All transition
functors fix objects, and $A$ is connected, so its object set
is exactly $C$. Given a word and an initial $c$, all intermediate
vertices of its path are forced by multiplication in $C$.
Every defining monoid relation can be used in any word context
at any such initial vertex, because the two substituted words
have the same product in $C$ and hence the same endpoint.
It follows in both directions that morphisms in the categorical
presentation are precisely elements $u\in C$ with $uc=c'$.
This proves (D.6.3); it does not assert that arbitrary categorical
colimits commute with nerves. The category $D_C$ has initial
object $1_C$, since the unique arrow $1_C\to c$ is labelled $c$.

\emph{Third, the homotopy colimit is simply connected.}
We spell out the fixed-object argument. For any diagram
$E:A\to\Cat$ with common object set $S$ and identity object
maps, the groupoid completion of $\int_A E$ has the arrows of
all $E(a)$ and, for each $f:a\to b$ and $s\in S$, a transport
arrow $t_{f,s}:(a,s)\to(b,s)$. The relations are composition
within the $E(a)$, composition of transports, and the
naturality squares for their arrows. Exactness implies totality
by the argument in D.2; hence $BA$ is nonempty and weakly
contractible. The groupoid completion $\Pi A$ is therefore
connected and thin. Choose an object $a_0\in A$. For each $s$,
conjugate every arrow to the objects $(a_0,s)$ along the unique
transport classes in $\Pi A$. Composition makes these choices
coherent, and all transport arrows become identities. The
remaining presentation consists exactly of the arrows and
relations of $\colim_A E$, with all arrows inverted. Thus
\[
 \Pi\!\left(\int_A E\right)
 \simeq\Pi\!\left(\colim_A^{\Cat}E\right).
 \tag{D.6.4}
\]
Apply this to $E=D$. By (D.6.3) its right side is the terminal
groupoid, so $B\int_A D$ is nonempty, connected and simply
connected.

Finally the homology spectral sequence for the Grothendieck
construction gives
\[
 E^2_{p,q}=T_p^A H_q(BD(a);\Z)
 \Longrightarrow H_{p+q}\!\left(B\int_A D;\Z\right).
\]
Exactness kills every term with $p>0$, and (D.6.2) kills every
term with $q\geq2$. Consequently the target has zero homology
in degrees at least two. Its $H_1$ vanishes by the just-proved
simple connectivity. The Hurewicz theorem, applied successively
to the CW realization of its nerve, makes it weakly contractible.
Quillen's Theorem~A now proves the assertion about $q_M$.
Combining with D.2 proves the final asphericity assertion for
$BC$.
\end{proof}

\begin{claim}{D.7}{Refuted}{Free monoid diagrams do not detect filteredness}
\label{D:not-filtered}
The implication asked in Part D.1 is false. The category
$\Delta_{\mathrm{inj}}$ is a nonempty connected nonfiltered
counterexample, and every diagram
$M:\Delta_{\mathrm{inj}}\to\mathrm{FreeMon}$ passes the specified
comparison with colimit in all monoids. In particular the
assertion in Part D.3 is true: every such monoid colimit has
aspherical classifying space. The same conclusion holds for
all the exact coloured-word categories $W_m$ in R6.
\end{claim}
\begin{proof}
We check R6 at the concrete step required. The normalized
integer cone in TOTASPH C.9 supplies, for every $[n]$, an
augmentation-one combination $t:[n]\to[N]$ for which $th$ is
independent of an injection $h:[k]\to[n]$. Its polynomial
verification uses gap monomials and the identity
$P_{t\delta^i}=y_iD_i(P_t)$; the coordinate-hyperplane
inclusion--exclusion construction makes all these expressions
equal and preserves value one at the all-ones point. Equalizing
cofaces implies equalizing all injections by their factorization
and the coface identities. Thus the affine category is filtered;
Isbell--Mitchell, Theorem~1, gives exactness. Their original
Theorem~3, read in the attached 1973 paper, states precisely this
filtered-affinization result, with its signed binary-tree cones.
For $W_m$, TOTEX E.2 lifts each $f:[l-1]\to[N-1]$ to the
colour-preserving map $i\mapsto mf(i-1)+w_i$ into
$(1,\ldots,m)^N$. For a coloured embedding $a:p\to w$, the
composite is $j\mapsto mf(a(j)-1)+p_j$, so the same affine
identities survive the lift. This verifies exactness of every
$W_m$ used here. Apply D.6. The two
point inclusions $[0]\rightrightarrows[1]$ cannot be equalized
by an increasing injection, proving nonfilteredness directly.
This counterexample is uniform over all free bases and all
word substitutions; it is not restricted to basis maps or to
the interval example.
\end{proof}

\paragraph{The positive-monoid issue for $\mathcal C_\delta$.}
D.4 answers Part D.4 by its expressly permitted alternative:
it gives a diagram of free groups with colimit in groups for
which the comparison fails. The group presentations (D.4.1)
include inverse letters. Abstract freeness and free-factor
inclusions do not specify compatible positive bases, and
forgetting the group structure gives monoids with nontrivial
units, not free monoids. Therefore D.4 does not by itself
give a diagram in $\mathrm{FreeMon}$. For a single relator one
can set $a=u_2u_1^{-1}=v_1v_3^{-1}$ and choose positive
expressions $u_2=au_1$, $v_1=av_3$. Across overlapping
hyperedges these choices impose simultaneous prefix equations;
the relative private-vertex argument does not prove their
compatibility. No assertion is made here that
$\mathcal C_\delta$ passes or fails all free-monoid tests.
The exact-category counterexample D.7 settles Part D.1
independently of that additional question.


\paragraph{Literature boundary.}
Fiedorowicz's Proposition~4.4 \cite{CDFiedorowicz1984} gives
another test for the conclusion about a monoid $C$: its
universal group-cover is the two-sided bar, so asphericity is
equivalent to vanishing of
$\Tor_i^{\Z C}(\Z,\Z[C^{\mathrm{gp}}])$ for $i>0$.
We read that proof. D.6 instead proves contractibility of the
comparison comma directly and therefore does not need an
amalgamation theorem or a flatness assertion for inclusions of
monoids. In particular it does not use the Lemma~5.2(a) of that
paper refuted by the attached Gray--Steinberg source. Searches
for exact colimits, free monoids, free categories, and
semicosimplicial monoids located no exact statement of D.6;
the source log records the opened texts and the unavailable
McDuff original. No claim of priority is made.
